\documentclass[11pt]{article}

% Shared notes settings for Elements of AIML lectures (UPES).
% Keep this file free of \documentclass to allow reuse via \input.

\usepackage[utf8]{inputenc}
\usepackage[T1]{fontenc}
\usepackage[a4paper,margin=1in]{geometry}
\usepackage{hyperref}
\usepackage{xurl}
\usepackage{booktabs}
\usepackage{amsmath}
\usepackage{amssymb}
\usepackage{listings}
\usepackage{xcolor}
\usepackage{enumitem}
\usepackage{parskip}

% Code listings (general-purpose).
\lstset{
  basicstyle=\ttfamily\small,
  keywordstyle=\color{blue},
  commentstyle=\color{gray},
  stringstyle=\color{teal},
  showstringspaces=false,
  columns=fullflexible,
  frame=single,
  framerule=0.2pt,
  breaklines=true
}

\setlist[itemize]{topsep=0.3em,itemsep=0.2em}
\setlist[enumerate]{topsep=0.3em,itemsep=0.2em}

% Simple per-section source line.
\newcommand{\notesources}[1]{\par\noindent{\small\color{gray}\textit{Sources:} \sloppy #1\par}}

% Unit-wise source strings for this subject (used by slides/notes).
% Path is relative to the lecture `latex/` folder (the compilation CWD).
\IfFileExists{../../../_shared/aiml-sources.tex}{%
  % Source strings for Elements of AIML (used by slides + notes).

\newcommand{\AIMLRepoURL}{https://github.com/tali7c/Elements-of-AIML}

\newcommand{\AIMLLabSources}{%
Provost and Fawcett, \textit{Data Science for Business}.;%
McKinney, \textit{Python for Data Analysis}.;%
WEKA Documentation (Waikato Environment for Knowledge Analysis), accessed 2026-02-28.;%
Matplotlib and Seaborn documentation, accessed 2026-02-28.%
}

\newcommand{\AIMLUnitISources}{%
Rich and Knight, \textit{Artificial Intelligence}, McGraw Hill, 2017.;%
Russell and Norvig, \textit{Artificial Intelligence: A Modern Approach} (4e), Ch. 1--2 (Introduction, Intelligent Agents).;%
Poole and Mackworth, \textit{Artificial Intelligence: Foundations of Computational Agents} (2e), selected sections.%
}

\newcommand{\AIMLUnitIISources}{%
Russell and Norvig, \textit{Artificial Intelligence: A Modern Approach} (4e), Ch. 7--9 (Logical Agents, First-Order Logic, Inference).;%
Huth and Ryan, \textit{Logic in Computer Science} (2e), selected sections on propositional and first-order logic.;%
Genesereth and Nilsson, \textit{Logical Foundations of Artificial Intelligence}, selected chapters.%
}

\newcommand{\AIMLUnitIIISources}{%
Mueller and Massaron, \textit{Machine Learning for Dummies}, For Dummies, 2016.;%
Mitchell, \textit{Machine Learning}, selected chapters on learning basics and evaluation.;%
Scikit-learn documentation, accessed 2026-02-28.%
}

\newcommand{\AIMLUnitIVSources}{%
Mueller and Massaron, \textit{Machine Learning for Dummies}, For Dummies, 2016.;%
James, Witten, Hastie, and Tibshirani, \textit{An Introduction to Statistical Learning} (2e), selected chapters.;%
Scikit-learn documentation, accessed 2026-02-28.%
}

\newcommand{\AIMLUnitVSources}{%
NIST, \textit{AI Risk Management Framework (AI RMF 1.0)}, 2023.;%
Barocas, Hardt, and Narayanan, \textit{Fairness and Machine Learning} (online book), selected sections.;%
Knaflic, \textit{Storytelling with Data}, selected chapters on communicating results.%
}
%
}{%
  % Faculty copies live under `private/` so their compile CWD is different.
  \IfFileExists{../../../../public/_shared/aiml-sources.tex}{%
    % Source strings for Elements of AIML (used by slides + notes).

\newcommand{\AIMLRepoURL}{https://github.com/tali7c/Elements-of-AIML}

\newcommand{\AIMLLabSources}{%
Provost and Fawcett, \textit{Data Science for Business}.;%
McKinney, \textit{Python for Data Analysis}.;%
WEKA Documentation (Waikato Environment for Knowledge Analysis), accessed 2026-02-28.;%
Matplotlib and Seaborn documentation, accessed 2026-02-28.%
}

\newcommand{\AIMLUnitISources}{%
Rich and Knight, \textit{Artificial Intelligence}, McGraw Hill, 2017.;%
Russell and Norvig, \textit{Artificial Intelligence: A Modern Approach} (4e), Ch. 1--2 (Introduction, Intelligent Agents).;%
Poole and Mackworth, \textit{Artificial Intelligence: Foundations of Computational Agents} (2e), selected sections.%
}

\newcommand{\AIMLUnitIISources}{%
Russell and Norvig, \textit{Artificial Intelligence: A Modern Approach} (4e), Ch. 7--9 (Logical Agents, First-Order Logic, Inference).;%
Huth and Ryan, \textit{Logic in Computer Science} (2e), selected sections on propositional and first-order logic.;%
Genesereth and Nilsson, \textit{Logical Foundations of Artificial Intelligence}, selected chapters.%
}

\newcommand{\AIMLUnitIIISources}{%
Mueller and Massaron, \textit{Machine Learning for Dummies}, For Dummies, 2016.;%
Mitchell, \textit{Machine Learning}, selected chapters on learning basics and evaluation.;%
Scikit-learn documentation, accessed 2026-02-28.%
}

\newcommand{\AIMLUnitIVSources}{%
Mueller and Massaron, \textit{Machine Learning for Dummies}, For Dummies, 2016.;%
James, Witten, Hastie, and Tibshirani, \textit{An Introduction to Statistical Learning} (2e), selected chapters.;%
Scikit-learn documentation, accessed 2026-02-28.%
}

\newcommand{\AIMLUnitVSources}{%
NIST, \textit{AI Risk Management Framework (AI RMF 1.0)}, 2023.;%
Barocas, Hardt, and Narayanan, \textit{Fairness and Machine Learning} (online book), selected sections.;%
Knaflic, \textit{Storytelling with Data}, selected chapters on communicating results.%
}
%
  }{%
    % Source strings for Elements of AIML (used by slides + notes).

\newcommand{\AIMLRepoURL}{https://github.com/tali7c/Elements-of-AIML}

\newcommand{\AIMLLabSources}{%
Provost and Fawcett, \textit{Data Science for Business}.;%
McKinney, \textit{Python for Data Analysis}.;%
WEKA Documentation (Waikato Environment for Knowledge Analysis), accessed 2026-02-28.;%
Matplotlib and Seaborn documentation, accessed 2026-02-28.%
}

\newcommand{\AIMLUnitISources}{%
Rich and Knight, \textit{Artificial Intelligence}, McGraw Hill, 2017.;%
Russell and Norvig, \textit{Artificial Intelligence: A Modern Approach} (4e), Ch. 1--2 (Introduction, Intelligent Agents).;%
Poole and Mackworth, \textit{Artificial Intelligence: Foundations of Computational Agents} (2e), selected sections.%
}

\newcommand{\AIMLUnitIISources}{%
Russell and Norvig, \textit{Artificial Intelligence: A Modern Approach} (4e), Ch. 7--9 (Logical Agents, First-Order Logic, Inference).;%
Huth and Ryan, \textit{Logic in Computer Science} (2e), selected sections on propositional and first-order logic.;%
Genesereth and Nilsson, \textit{Logical Foundations of Artificial Intelligence}, selected chapters.%
}

\newcommand{\AIMLUnitIIISources}{%
Mueller and Massaron, \textit{Machine Learning for Dummies}, For Dummies, 2016.;%
Mitchell, \textit{Machine Learning}, selected chapters on learning basics and evaluation.;%
Scikit-learn documentation, accessed 2026-02-28.%
}

\newcommand{\AIMLUnitIVSources}{%
Mueller and Massaron, \textit{Machine Learning for Dummies}, For Dummies, 2016.;%
James, Witten, Hastie, and Tibshirani, \textit{An Introduction to Statistical Learning} (2e), selected chapters.;%
Scikit-learn documentation, accessed 2026-02-28.%
}

\newcommand{\AIMLUnitVSources}{%
NIST, \textit{AI Risk Management Framework (AI RMF 1.0)}, 2023.;%
Barocas, Hardt, and Narayanan, \textit{Fairness and Machine Learning} (online book), selected sections.;%
Knaflic, \textit{Storytelling with Data}, selected chapters on communicating results.%
}
%
  }%
}


\title{Elements of AIML Lab\\Experiment 08 (After-submission Reference): PCA (Dimensionality Reduction)}
\author{Tofik Ali}
\date{\today}

\begin{document}
\maketitle
\tableofcontents

\section{How to use this reference}
This reference is meant to be shared \textbf{after you submit} Experiment~08.
It clarifies:
\begin{itemize}[leftmargin=*]
  \item why scaling is needed before PCA,
  \item how to choose number of components using explained variance,
  \item how to avoid leakage when PCA is used before a model.
\end{itemize}

\section{PCA: what it learns}
PCA finds new axes (principal components) such that:
\begin{itemize}[leftmargin=*]
  \item PC1 captures maximum variance,
  \item PC2 captures maximum remaining variance subject to being orthogonal to PC1,
  \item and so on.
\end{itemize}
Each principal component is a weighted combination of original features.

\section{Worked example (end-to-end)}
\subsection{Step 1: load data and select features}
\begin{lstlisting}[language=Python]
import pandas as pd

df = pd.read_csv("processed.csv")  # or your dataset
X = df.select_dtypes("number")     # PCA requires numeric features
\end{lstlisting}
If you have a label column, keep it separately (do not include in PCA features).

\subsection{Step 2: scale + fit PCA}
\begin{lstlisting}[language=Python]
import numpy as np
from sklearn.preprocessing import StandardScaler
from sklearn.decomposition import PCA

scaler = StandardScaler()
X_scaled = scaler.fit_transform(X)

pca = PCA()
Z = pca.fit_transform(X_scaled)

evr = pca.explained_variance_ratio_
cum = np.cumsum(evr)
print("First 10 EVR:", evr[:10])
print("First 10 cumulative:", cum[:10])
\end{lstlisting}

\subsection{Step 3: pick number of components}
Two common choices:
\begin{itemize}[leftmargin=*]
  \item choose the smallest $k$ such that cumulative EVR $\ge 0.90$ (90\% variance),
  \item or choose $k$ at the ``elbow'' of the cumulative curve (diminishing returns).
\end{itemize}
You can also let scikit-learn choose by ratio:
\begin{lstlisting}[language=Python]
pca90 = PCA(n_components=0.90)
Z90 = pca90.fit_transform(X_scaled)
print("Chosen components:", pca90.n_components_)
\end{lstlisting}

\subsection{Step 4: plot cumulative explained variance}
\begin{lstlisting}[language=Python]
import matplotlib.pyplot as plt

plt.plot(cum, marker="o")
plt.axhline(0.90, color="red", linestyle="--", label="90%")
plt.xlabel("Number of components")
plt.ylabel("Cumulative explained variance")
plt.grid(True)
plt.legend()
plt.show()
\end{lstlisting}

\subsection{Step 5: interpret components (loadings)}
Loadings show which original features contribute to each PC.
\begin{lstlisting}[language=Python]
import pandas as pd

loadings = pd.DataFrame(
    pca.components_.T,
    index=X.columns,
    columns=[f"PC{i+1}" for i in range(pca.n_components_)]
)

print(loadings["PC1"].sort_values(key=abs, ascending=False).head(8))
\end{lstlisting}
\textbf{Interpretation tip:} large magnitude loading $\Rightarrow$ that feature strongly influences that component.

\section{Leakage-safe PCA (important)}
If you use PCA before training a model:
\begin{itemize}[leftmargin=*]
  \item split train/test first,
  \item fit scaler+PCA on training only,
  \item transform test using fitted objects,
  \item ideally use a pipeline.
\end{itemize}
Example:
\begin{lstlisting}[language=Python]
from sklearn.pipeline import Pipeline
from sklearn.linear_model import LogisticRegression

pipe = Pipeline([
    ("scaler", StandardScaler()),
    ("pca", PCA(n_components=0.90)),
    ("clf", LogisticRegression(max_iter=1000)),
])
\end{lstlisting}

\section{Troubleshooting}
\begin{itemize}[leftmargin=*]
  \item \textbf{PCA fails:} check for non-numeric columns and missing values.
  \item \textbf{EVR is low:} your features may be weakly correlated; PCA may not compress much.
  \item \textbf{PC plot looks noisy:} PCA is not a classifier; separation is not guaranteed.
\end{itemize}

\notesources{\AIMLLabSources}
\end{document}

